\begin{abstract}
Simulation studies often assess virtual control arms by comparing estimated
effects with a requested treatment effect. A recovery curve can look well
calibrated even when the estimate exactly reproduces a reference statistic of
the simulated sample. We propose a simple audit: pre-specify that realised
reference and report the estimate's residual against it alongside error against
the requested parameter. In a single-cell oracle analysis, the residual was
zero across 1.7 million replicates, yet direct calibration did not support the
cell-count threshold previously justified by the recovery curve. A stratified
patient simulator shows the same equality for saturated G-computation,
empirical-propensity weighting and non-cross-fitted augmented weighting.
The result depends on both the reference and the trial design: exact balance
within strata makes even the unadjusted contrast reproduce the reference,
while censoring separates observed-data and latent-time comparisons.
A zero residual flags possible functional reuse, not an invalid estimator; a
nonzero residual does not establish validity. A residual-to-generator-error
ratio quantifies departures from equality without attributing fractions of
total error. The audit clarifies what a recovery curve measures and why interval
coverage and abstention rules require separate evaluation.
\end{abstract}
